% The duty cycle as a state machine, with the marker code each state asserts.
\begin{tikzpicture}[font=\sffamily\scriptsize, node distance=9mm and 16mm]

\node[initial] (i) at (0,0) {};
\node[state, text width=20mm, right=7mm of i] (sleep) {SLEEP\\{\tiny code 000}};
\node[state, text width=20mm, right=of sleep] (wake) {WAKE\\{\tiny code 001}};
\node[state, text width=20mm, right=of wake] (sense) {SENSE\\{\tiny code 010}};
\node[state, text width=20mm, below=16mm of sense] (comp) {COMPUTE\\{\tiny code 011}};
\node[state, text width=20mm, left=of comp] (send) {SEND\\{\tiny code 100}};
\node[state, fill=hwcol, text width=20mm, left=of send] (fault) {FAULT\\{\tiny code 111}};
\node[final, left=10mm of fault] (f) {};

\draw[trans] (i) -- (sleep);
\draw[trans] (sleep) -- node[above]{RTC} (wake);
\draw[trans] (wake) -- node[above]{clocks up} (sense);
\draw[trans] (sense) -- node[right]{samples ready} (comp);
\draw[trans] (comp) -- node[above]{feature done} (send);
\draw[trans] (send) -- node[above]{no reply} (fault);
\draw[trans] (fault) -- node[above]{halt} (f);
\draw[trans] (send.north) -- ++(0,6mm) -| node[lbl, pos=0.25, above]{stub returns} (sleep.south);

\node[umlnote, text width=44mm, anchor=north west] at (-0.6,-4.2)
  {One store per transition. The code is raised on entry and never cleared
   inside a state, so every instant of the cycle carries exactly one code and no
   charge falls between two phases.};
\node[umlnote, text width=44mm, anchor=north west] at (4.8,-4.2)
  {Code 000 is also the reset state, so an unpowered or halted board decodes as
   sleep rather than as an undefined phase. Code 111 is raised by the fault
   handler and never cleared.};
\end{tikzpicture}
